\scanpage{289}
% archive sheet number: 290
On considère les foncteurs d'inclusion
\[
i_N:\mathcal M_N\hookrightarrow\mathcal M\qquad i_{\deleted{\unsure{N}}}^{N}:\mathcal M^{N}\mathrel{\unsure{-|\to}}\mathcal M\tag*{(9)}
\]
plus généralement, pour $N\leq N'$
\[
i_{N,N'}:\mathcal M_{N'}\hookrightarrow\mathcal M_N\qquad i^{N',N}:\mathcal M^{N}\mathrel{\unsure{\subset|\hookrightarrow}}\mathcal M^{N'}\tag*{(9.1)}
\]

\deleted{On va étudier les}

donc on a
\[
i_N\,i_{N,N'}=i_{N'\deleted{\unsure{,N}}}\quad,\quad i^{N'}i^{N',N}=i^{N}\tag*{(9.2)}
\]
\deleted{plus \unsure{g}} et si $N\leq N'\leq N''$, d'où \deleted{$\unsure{i_N}$}
\[
\left\{\begin{aligned}
\mathcal M_{N''}\subset\mathcal M_{N'}\subset\mathcal M_N\ \deleted{\text{\unsure{et}}}&\qquad\mathcal M^{N}\subset\mathcal M^{N'}\subset\mathcal M^{N''}\\
i_{N,N''}=i_{N,N'}i_{N',N''}&\qquad i^{N'',N}=i^{N'',N'}i^{\unsure{N'N|N',N}}
\end{aligned}\right.\tag*{(9.3)}
\]

On va étudier la suite de foncteurs adjoints associés à $i_N,i_{N,N'}$ d'une part, $i^N,i^{N,N'}$ d'autre part.

\uline{Proposition} 1 \quad $i_N$\inserted{$:\mathcal M_N\hookrightarrow\mathcal M$} s'insère dans une suite de 4 foncteurs adjoints
\[
q_N\unsure{,}\quad i_N\quad \tau_N\quad \varphi_N\tag*{(10)}
\]
(chaque foncteur \deleted{\unsure{étant|adjoint}} parmi $i_N,\tau_N,\varphi_N$ étant adjoint à droite de celui écrit à \unsure{sa|2} gauche) définis ainsi :

a) $\tau_N:\mathcal M\longrightarrow\mathcal M_N$ (foncteur \unsure{tronquage|troncature} à droite)
\[
\tau_{\unsure{N-1|N}}(P)\mathrel{\underset{\substack{\downarrow\\\text{hom d'adjonction}}}{\subset}}\deleted{\unsure{i}}P
\quad\text{(sous objet de $P$)}\quad.
\]
\[
\begin{aligned}
&\overset{\Vert}{i_N\tau_N(P)}\\
(\tau_N(P))_{i}&=\left\{\begin{aligned}
P_i&\quad\unsure{=}\ \text{si }i\geq N\\
0&\quad\text{si }i<N
\end{aligned}\right.
\end{aligned}\tag*{(10.1)}
\]
